\section{A common short-range covariance on every period}
\label{sec:actual-reader-short-covariance}

This section records the finite-period covariance output used by the local
reader comparison. The construction keeps one bed, one forest, one blocking
history, and one cover row throughout. In particular, the row is not chosen
again when the period changes. The exact curvature covariance on the tree
slice is identified in Section~\ref{sec:actual-reader-tree-covariance}; we
use that identity here and construct its common exponentially localized
cover representative and its finite-range approximation.

\subsection{Period folding and the common inverse}

Let $\mathcal I$ be the finite set of asked coordinates for a
fixed forest and gauge branch. Its labels include the intra-block residue
and colour. Let $\mathcal C$ denote the curvature labels, likewise
including the intra-block cochain type and colour. The lattice coordinate
in a cover row is always a coarse cell base, not a fine site. Write $d_n$
for shortest $\ell^1$ distance between coarse bases on the period-$n$
torus. A cover row is a family $P_{ab}(z)$, $a,b\in\mathcal I$,
$z\in\mathbb Z^4$. Its period-$n$ matrix
is the absolutely convergent periodization
\[
 P^{(n)}_{(x,a),(y,b)}
   =\sum_{z:\ y=x+z\;\mathrm{mod}\;n}P_{ab}(z).
\]
A bound $\sum_z e^{s|z|_1}|P_{ab}(z)|\le M$ makes this sum absolute and
bounds each row uniformly in the period. All sums below use the same
coordinate and colour labels.

\begin{lemma}[Folding a coercive row to every period]
\label{lem:ex-reader-folding}
Suppose $P$ is a real symmetric cover row, has an exponential row bound
of width $s>0$, and there are $a>0$ and $n_0$ such that
\[
 \langle v,P^{(n)}v\rangle\ge a\|v\|_2^2
 \quad\text{for every }n\ge n_0.
\]
If the periodizations commute with pullback under every quotient map
$\mathbb T_m^4\to\mathbb T_n^4$ with $n\mid m$, then the same lower bound
holds for every positive period.
\end{lemma}
\begin{proof}
Fix $n$. Choose $m=(n_0+1)n$ and let $g$ be coordinatewise reduction
modulo $n$. Then $g$ is surjective and commutes with lattice translations.
Pull a vector $v$ on period $n$ back to $v\circ g$. Each coordinate of
$v$ has exactly $|\ker g|$ preimages, so
\[
 \|v\circ g\|_2^2=|\ker g|\|v\|_2^2.
\]
At a coordinate $(x,i)$, the left side is
\[
 \sum_{j\in\mathcal I}\sum_{z\in\mathbb Z^4}
       P_{ij}(z)\,v(g(x+z),j).
\]
Since $g(x+z)=g(x)+g(z)$, group the absolutely convergent sum by the
residue $g(z)$. The result is $(P^{(n)}v)(g(x),i)$, so
$P^{(m)}(v\circ g)=(P^{(n)}v)\circ g$. Taking inner products and
using the same fiber count on both sides gives
\[
 |\ker g|\langle v,P^{(n)}v\rangle
 \ge a|\ker g|\|v\|_2^2.
\]
The fiber size is positive, so division proves the claim. This argument
also applies with the finite colour coordinate present: the quotient acts
only on the spatial coordinate.
\end{proof}

The inverse row is obtained without a Fourier-symbol assumption. For each
period, take the ordinary matrix inverse of $P^{(n)}$. A uniform positive
lower bound gives the finite-volume resolvent bound. The exponential row
estimate gives a uniform weighted row and column bound on every
periodization: the internal label set is finite and each periodic
 displacement sum is bounded by the corresponding cover sum. The diagonal
conjugation estimate of Section~\ref{sec:local-covariance}, applied to
these finite matrices, gives constants $C_1,r_1>0$, independent of the
period, such that each inverse entry is bounded by
$C_1e^{-r_1d_n(x,y)}$. This is the finite-dimensional weighted-resolvent
estimate used in the period-independent inverse construction; no spectral
lower bound at complex momentum is assumed.
For every factorial side $N=k!$, choose one centered shortest
representative in each residue class, coordinatewise in
$[-\lfloor N/2\rfloor,\lceil N/2\rceil)$. Define the cover coefficient
read from its inverse at $(0,a)$ and $(z,b)$ when $z$ is this representative,
and set it to zero otherwise. For every retained coefficient,
$d_N(0,z)=|z|_1$, so the bound $C_1e^{-r_1|z|_1}$ holds uniformly in both
$z$ and $k$. For every fixed $z$, it is retained for all sufficiently large
$k$. Enumerating the countable set of triples and taking successive
subsequences, followed by the diagonal subsequence, gives simultaneous
coefficient convergence. Passing to finite
partial sums retains the exponential bound. At any fixed smaller width,
the remaining tails are dominated by a constant times
$e^{-c|z|_1}$, a summable function on $\mathbb Z^4$; hence the inverse rows
converge in the weighted row norm.

The limit has the required inverse on every period, not merely on the
factorial sequence. Indeed, for $n\mid m$, pullback intertwines the two
periodic convolution matrices. Since both are invertible, it also
intertwines their inverses: apply the forward intertwining identity and
uniqueness of the solution to the finite linear system. Periodizing the
cover limit along any fixed period is justified by the common exponential
majorant. Passing to the limit in the factorial periods divisible by $n$
therefore identifies that periodization with $(P^{(n)})^{-1}$. This is the
period-independent inverse construction used below.

\subsection{The selected tree slice on the same history}

Fix a blocking patch, gauge, kernel, forest, and bed $H$. At each period
$n$, repeat the same forest by the zero-offset lift. Let $\mathcal A_n$
be the fine one-form coordinates and let $\mathcal I_n$ be the active
asked-coordinate set. The cover coefficients $P_{ij}(z)$ are the common
full, uncut row on the finite asked-label set $\mathcal I$. Periodization
sums these coefficients over the displacement classes and gives exactly
the restriction of the finite tree precision to the repeated forest.

The local construction has this concrete input. Fix one odd block side
$L$, one patch and gauge branch, and one forest on the reference three-side
block. Repeat that forest at every period by the same zero-offset rule.
Fix a bed $H$, a clip, and constants $s,M,a,V,K_{\rm pen}$ with
$s,a,K_{\rm pen}>0$, $M\ge0$, and $V>0$. Let $G_n$ and $\Delta_n$ be its metric and diagonal matrices
on period $n$. Assume $H$ has strip width $s$ and bound $M$, and that at
each finite period these matrices are symmetric and satisfy
\[
 a\|u\|_2^2\le \langle u,G_nu\rangle\le V\|u\|_2^2,
 \qquad
 a\|u\|_2^2\le \langle u,\Delta_nu\rangle\le V\|u\|_2^2.
\]
Also assume the exact root-gradient identity
\eqref{eq:reader-root-gradient}, the same-bed stationarity identities
\eqref{eq:reader-stationarity-input}, and the two-form transport identity
\eqref{eq:reader-curvature-intertwiner} from
Section~\ref{sec:actual-reader-tree-covariance}. The displayed numerical
bounds produce the gap; these algebraic identities identify the
conditional covariance. They are distinct inputs. The construction fixes
its positive width $t$ and row constant $A$ from the patch, kernel, branch,
forest template, and $s,M,a$ before selecting the bed within this class;
the same pair works uniformly for every bed satisfying the stated strip
and floor bounds.

The finite block estimate of Section~\ref{sec:path-attachment-producer} gives the
lower bound
\[
 \frac{\min(a,K_{\rm pen})}{1+\beta+\gamma},
 \qquad \beta,\gamma\ge0,
\]
for every coarse side at least $100R+3$, where $R$ is the macro radius
and $K_{\rm pen}$ is the prescribed penalty. Here $X$ is the fixed
patch-offset set and $C_{\rm edge}$ is the edge-comparison constant. The
finite costs are
\[
 \beta_\ell=(2C_{\rm edge}^2)|\operatorname{Site}(L)|\,|\operatorname{Axis}|
 (|X|\,|\operatorname{PlaqAxes}|)\,|X|,
 \qquad \gamma_\ell=2|\operatorname{Site}(L)|.
\]
They are the curl and rooted-probe coefficients proved in
Section~\ref{sec:path-attachment-producer}; in the gap formula above
$\beta=\beta_\ell$ and $\gamma=\gamma_\ell$.
The cover coefficients are obtained by composing the bed row with the
fixed local blocking maps. The weighted row norm of a product is bounded
by the product of the weighted row norms, since
$|z+w|_1\le |z|_1+|w|_1$. The blocking maps have finite coefficient sums,
so the resulting exponential bound is finite and independent of the
period. Periodizing these same coefficients gives the repeated full row,
whose restriction is the tree precision. Thus it has the displayed gap
at all sufficiently large periods. For any requested period $n$, take
$m=(100R+4)n$, reduce modulo $n$, and pull a vector back. The quotient is
surjective and every fiber has the same positive cardinality. Directly
reindexing the cover convolution gives
$P^{(m)}(v\circ g)=(P^{(n)}v)\circ g$, while both the quadratic form
and squared norm of $v\circ g$ are multiplied by $|\ker g|$. Dividing
transfers the same gap to period $n$.

Apply the weighted finite-resolvent estimate to this common row. Its
inverse is one exponentially localized cover row whose period action is
the inverse tree precision at every period. Let $\mathcal R$ be the fixed
rooted path-measurement matrix from asked-coordinate vectors to curvature
coordinates, so its rows are indexed by $\mathcal C$ and columns by
$\mathcal I$. The rooted gauge projection
and the identity between the root correction and the gradient evaluation
show that the curvature is unchanged by removing forest coordinates.
Thus the matrix
\[
             Z=\mathcal R\,P^{-1}_{\mathrm{row}}\,\mathcal R^{\mathsf T}
\]
has, on each period, the same restriction to active curvature coordinates
as the finite tree-slice covariance for every admissible coarse period
$n\ge3$. More explicitly, the matrix on the
slice is $(DI_F)A_F^{-1}(DI_F)^{\mathsf T}$, which equals
$P_f-JP_cJ^{\mathsf T}$ by \eqref{eq:reader-tree-conditional-covariance}.
This is an exact equality. Its proof
uses the same bed and the same forest on both sides; it does not replace a
curved bed by the flat one or compare two different histories. The
finite-dimensional curl-through identity and the resulting exact
conditional curvature formula are the ones proved in
Section~\ref{sec:actual-reader-tree-covariance}.

The row norm of $Z$ is bounded by the product of the row norms of the two
rooted path rows and the inverse row. For convolution, the weighted row
norm is submultiplicative: expand the product, use
$|z+w|_1\le |z|_1+|w|_1$, and sum the resulting nonnegative majorant.
Therefore $Z$ has a positive exponential width, with a finite constant
fixed before the period is chosen.

If the desired period matrices are symmetric but this particular cover
representative is not symmetric term by term, replace it by
\[
 Z^{\mathrm{sym}}_{ab}(z)
 =\tfrac12\bigl(Z_{ab}(z)+Z_{ba}(-z)\bigr).
\]
Its period matrix is the average of the original matrix and its transpose.
Since the exact tree covariance is symmetric, this replacement leaves every
period matrix unchanged. The triangle inequality shows that its exponential
row bound is no larger than the original bound. The reversal is essential:
transposition reverses the displacement as well as exchanging the two
indices.

\subsection{Cutting to short paths}

For a scale parameter $x>1$, set the integer radius
$B_x=\lceil(\log x)^2\rceil$ and the natural-number cut radii
$h_x=\lfloor B_x/10000\rfloor$ and $H_x=\lfloor B_x/8\rfloor$.
Their real uses below mean their natural casts. Choose a symmetric
coarse-cell displacement cut which retains every coefficient with
$|z|_1\le h_x$ and retains none with $|z|_1>H_x$. Write $Z^{\le x}$ for
the cut row. If
$\sum_z e^{t|z|_1}|Z_{ab}(z)|\le A$, then for every finite-period entry the triangle inequality and the tail
estimate give
\[
 \left|Z^{(n)}_{(u,a),(v,b)}
       -(Z^{\le x})^{(n)}_{(u,a),(v,b)}\right|
 \le A e^{-t h_x}.
\]
The same estimate holds with row and column interchanged. The cut matrix
is symmetric because the retention rule is invariant under
$(a,b,z)\mapsto(b,a,-z)$. Also, if the separation of the two spatial
coarse parent-base addresses (the addresses obtained by splitting the fine
site into its coarse base and intra-block residue) is greater than $H_x$,
its cut period-matrix entry is zero: every retained coarse displacement
has length at most $H_x$, so none can connect those bases modulo the period. These two facts are the only geometric
properties of the cut used in the short-history comparison.

In the source normalization the row-tail payment for all output labels is
\[
 \operatorname{cheque}(t,A,x)
 = |\mathcal C|A\exp(-t h_x).
\]
For the installed two-form path row, the incoming cover bound has a
specific source. In cover-row notation it is
\[
 \Psi=D_f^{-1}\,\mathrm{Fine}(T)\,D_c,
\]
where $D_f$ is the fine diagonal, $D_c$ the coarse diagonal, and
$\mathrm{Fine}(T)$ is the installed tower map. Assume the corresponding
same-history strip receipts: $D_f^{-1}$ has bound $M_f$, $D_c$ has bound
$M_c$, and the moving tower row has bound $M_T/L^2$. The weighted
row-product estimate gives
\[
 \Psi\text{ has strip bound }
 \frac{m_2^2 M_fM_T M_c}{L^2}.
\]
In the selected case where the three incoming factors are all bounded by
$M_{\rm in}$, put $\psi_{\rm budget}(u)=m_2^2u^3$; then the installed row
has bound $\psi_{\rm budget}(M_{\rm in})/L^2$. The factor $L^{-2}$ is
already in the moving tower receipt, not produced by installation;
installation preserves its per-entry bound and summing over output labels
supplies the cheque factor. The paid movement error is therefore
$\operatorname{cheque}(s,\psi_{\rm budget}(M_{\rm in})/L^2,x)$. This is an
incoming budget, separate from the turn-row cheque.

Define the short transport matrix on a period by keeping each installed
transport entry when the two coarse parent bases have separation at most
$H_x$, and setting it to zero otherwise. The cut path row has zero at all
omitted addresses. The installed-row estimate therefore bounds, for each
input coordinate, the omitted row sum by
$\operatorname{cheque}(s,\psi_{\rm budget}(M_{\rm in})/L^2,x)$; in
particular every omitted entry has this bound, while retained entries are
unchanged. This is the same two-case finite-range comparison as for
$Q_n^{\rm short}$ below.

\begin{proposition}[Cutting the constructed same-history covariance]
\label{prop:ex-reader-short-covariance}
Under the fixed branch, repeated-forest, bed strip/floor, and same-bed
factorization hypotheses above, construct $Z$ by the common inverse and
rooted measurement. If a symmetric cut retains every displacement through
$h_x$, has support through $H_x$, and its row-tail payment is at most
$\varepsilon$ in both orientations, then $Z$ has the exact finite
tree-slice curvature covariance at every admissible coarse period $n\ge3$,
while the cut differs entrywise from that covariance by at most
$\varepsilon$. Its entries vanish at separations greater than $H_x$.
\end{proposition}
\begin{proof}
The finite tree bound gives the gap above the stated side threshold, and
Lemma~\ref{lem:ex-reader-folding} transfers it to every period. The
period-independent inverse and rooted measurement produce $Z$ with the
exact period matrices for the admissible coarse periods $n\ge3$.
Symmetrize by displacement reversal; exact symmetry of the period
covariances preserves every period output. The cut-tail
estimate gives the entrywise error, and the support-radius condition gives
the vanishing statement.
\end{proof}

There is a useful finite-range version of the cut output. On period $n$,
let $C_n=Z^{(n)}$ be the exact covariance and define
\[
 Q^{\rm short}_{n}(u,a;v,b)=
 \begin{cases}
 (C_n)_{(u,a),(v,b)},&d_n(u,v)\le H_x,\\
 0,&d_n(u,v)>H_x.
 \end{cases}
\]
This is symmetric and has range $H_x$. If $d_n(u,v)\le H_x$, it agrees
with the covariance by definition. If $d_n(u,v)>H_x$, the cut period
matrix vanishes by the support condition, so the bound from
Proposition~\ref{prop:ex-reader-short-covariance} gives
\[
 |(C_n)_{(u,a),(v,b)}-Q^{\rm short}_{n}(u,a;v,b)|
 \le \operatorname{cheque}(t,A,x).
\]
Thus the short comparison used at a reader is exact on retained pairs;
the only error is the omitted covariance tail. This is the two-case form
needed for a finite-range covariance comparator on one fixed history.

\subsection{How the estimate enters the path telescope}

For clarity, the quantitative input used by the path telescope consists of
two payments, not a free approximation assertion. If $\Psi$ is the
installed two-form path row, let $m_2$ be the number of two-form labels,
set $\psi_{\rm budget}(u)=m_2^2u^3$, and let $M_{\rm in}$ be its
pre-installation bound. The movement payment is
\[
 m_2\,\frac{\psi_{\rm budget}(M_{\rm in})}{L^2}
       \exp(-s h_x)\le \varepsilon.
\]
The second payment asks for a same-bed row $Z$, a width $t>0$, and a
bound $A>0$ such that the preceding cut estimate is at most
$\varepsilon$. The path part of the telescope uses the movement receipt
through the triangle inequality and the finite row sum of the installed
path data; the covariance part uses the turn receipt and the exact
same-bed transfer. These are separate clauses. This section proves the
row and covariance estimate, but it does not derive either price from the
scale schedule.

Suppose two consecutive levels of the selected history have the same
blocking clock. In the covariance clause, truncate the rooted row at
radius $H_x$. For endpoints farther apart, the truncated matrix entry
vanishes. For the remaining pairs, the exact same-bed transfer identifies
the uncut entry with the finite tree conditional covariance, and the cut
estimate bounds its replacement error. In the path clause, the installed
path row is handled by its finite row-sum estimate and the separate
movement receipt. Both clauses retain the original stamped history. The
selected schedule must allocate the movement and turn prices to their
respective requested tolerances; the field-path bound and capital budget
remain additional hypotheses of the telescope.

The large-period gap, the strip width, and both payments are properties
of the selected bed and the selected history. Their production is not
proved by the finite covariance algebra above. In particular, this result
must be combined with the separate schedule and paid-budget construction
before it can be used as an unconditional Yang--Mills construction step.
